\section{Ascesis and Anti-Europe}

\begin{quotex}
This is authored by \textbf{Havismat} (\textbf{Guido De Giorgio}), from Volume 2 of Introduction to Magic. We see more of the paradoxical style of his writing. He reveals some personal details about his time in Tunisia and his experience with Sufis. See Short Note on Woman in East and West\footnote{\url{https://gornahoor.net/?p=6420}} for an interesting note that was in the early editions, but not the final edition of Introduction to Magic.
\end{quotex}

Regarding expression. There is the traditional one, the doctrinal one, the personal one. The first is integral, complete because of a pure intuitive order (Upanishad): the second is rigorous and complete, but presupposes the first: the third is poetic, strong, \emph{poignant} --- in the Sufi, for example --- these three forms are totally complete and definitive in spite of the enormous differences: in one word, orthodox, in the etymologically rigorous and therefore true meaning of the term. These three forms of expression are superior in adhering to the inexpressible, eliminating everything that, as individual, alters and makes this same adherence opaque. These three forms of expression are absolutely innocuous: for those who neither know nor intuit, they are like water between the fingers; for those who intuit, they are cast, fleeting and revealing; for those who know, they are bends of teaching, they are support, discoveries and, in ascesis, stellar vertices.

Beyond all the hesitations, efforts, breakthroughs, perpendicular invasions, or serpentine vortices of personal labors of resolution and culmination, the metaphysical domain of the great Tradition shows an absolute, very clear transparency. In everything, one can deceive and be deceived---and the more one knows, the more one realizes that --- except on that level that can be intuited, through convergences, or ``integrated'', but does not cease being what it is: a point, a knot of points of continuous connection that remains an unfathomable and intangible heaven, but only heaven, for those who intuit, that-cannot-be-said for those who know.

If poetry is the marginal game that opens the cracks of the Illusion, so the Poet, who opens and closes his eyes, is to be preferred to a minor expert and to a false mage; he is true --- flashing, cultured, and lost depths --- but always on the true axis, on the axis of the truth. This is why a tradition (in the integral sense) admits and justifies, without disdain, the one who knows and who does not know as long as the \emph{axis is unique}, as long as there is no prevarication, as long as the one who does not know or who only intuits says: ``It is here, it is absolutely here, also and \emph{especially} (and this is the enormity of the Orient) if I will not have it and blessed are those that have it and who proceed.''

India has achieved that integrally, extraordinarily: I refer in a special way to the Dharmanga and Acrama, that is, to the integral tradition in the totality of its accomplishments and the resolution of the accomplishments. The West does not know that: there can be no true resolution, ``asceticism'', accomplishment, if Illusion is not placed in its own domain, i.e., if its members are not recomposed in a way that everything converges in that which is the Root of Illusion.

Elite, that is good: they are provisional means. The martyrs of the West have not moved the overturning of the West by a millimeter: the tradition, through a thin thread, of the West, has not subverted this cadaver that plays at resurrection. An elite is not the tradition: an elite is a vein, a precious vein, but a vein: it wants the boulder, they want other veins, and it is necessary that all the veins converge, although the central one alone, the most hidden, it goes regally and, hiding itself, submerging itself, dominates.

If there is no space, there is no echo, if there is no silence, there is no Voice, if there is no unanimity, there is no song. To create adaptations, and what is begin done in the West if not this, and for centuries? It is necessary that things sing, the boulders, the metals, the woods, the thoughts, the stars. Everything must sing, from the blade of grass to Brahman.

But in the West one plays masquerades: a mask appears and disappears: a thrill of passion --- ``he is a genie'' --- seen, known, buried: and he draws away: and by degrees the West submerges itself. Do you think of the price of the thing, of the object: what is one thing worth now? Nothing. A bucket, a book, a statue, a house, a city, the world: what is it worth? Nothing. What is it to give birth? Nothing. What is it to die? Nothing. What is it to think? Nothing.

Man now steals life: if the hunt is in his pocket, he touches it fearfully: he escapes from it. Here you are. He lives for nothing, dies for nothing. There is no sense. He plays with ghosts: the Lemurian age: individual, humanity, society, family, science. Those do not exist. That is outside Tradition: Tradition does not know ghosts: that is why it is not only about showing and suggesting a unique higher type, but of establishing the conditions for which a tradition, which includes all the variously ordered activities but along a single axis, can exist. To hurl thrills of absolute liberation to men who do not even know how basic human existence should be established in the field of action, but always related to a center that is not its essence, is dangerous and absurd. \emph{Either life is a rite} or it is nothing: either everything reacquires a symbolic character, under the general type of the offering, or else nothing remains.

The \emph{tragic sense of life}, for which modern men have so much sympathy, is also the element of dissolution and above all irresolution. Life is not tragic: tragedy has created the desperate man. Life is that which is and is that which must be. Pessimism harms more than optimism. The modern idea of ``value'' is already a construction like that of \emph{will} and \emph{power} in another field. Neither optimism nor pessimism, but \emph{to give meaning to action}, to direct it on the basis of a boundary of resolution, to bend it to give everything down to the foundation; in everyone, not in the few. The few will be here: but there will always be those-who-go, not the grim despisers, but those who watch very calmly the turmoil of the waters from above.

Today one ``builds'' on the whole line: such tragic sense sweeps things away, it is the Rembrandtian sense of the shadow that sweeps away the truth: the white, the neutral, it is a key, the psychological key that slaughters the very placid grass of the meadow.

Ascesis is not asceticism. Asceticism is a deviation of ascesis. To transcend means to surpass that which is called reality and to see this reality under a different aspect: one cannot achieve the one without automatically getting the other. Therefore, whoever transcends, changes his eyes and by changing his eyes, everything is changed. The relationship Ascesis-Reality is the true one: supposing that one realizes Fulfillment, i.e., that individuality is re-solved, reality is play, i.e., tentacles that now open up, now disappear according to whether the eye opens or closes. The concept of Ascesis is precisely this: the elimination of everything that makes of the I that which it is: one denies, one burns, one dissolves up to that, at the boundary of the end, the resolution is completed. Then this is absolutely exact: ``\emph{Quicumque Deum intelligit Deus fit}'' ``whoever knows God becomes God'', from the \emph{Mundaka Upanishad} , but that \emph{quicumque} is no longer \emph{quicumque} when the transition from \emph{intelligere} to \emph{fieri} is complete. but that ``who'' is no longer the same ``who'' when the transition from ``knowing'' to ``becoming'' is complete. And I would like to say: \emph{Secretum tegendum, secretum tegendum} literally ``secret to be kept'', the Latin translation of \emph{Upanishad}, now more than ever. I know of someone who went to a Sufi and demanded from him like a European: in French, in the original ``Initiate me into Sufism.'' The Sufi looked at me, I smiled, and he responded: ``But my dear sir, if you are not yourself a Sufi, how do you count on becoming one?'' The European said: ``But, in the end, through initiation?'' ``Initiation!'' said the Sufi watching the sea with eyes grinning. And I answered for him roughly: ``There is no initiation: either beyond you, or nothing.''

This man, who is still looking, never finds, in spite of the lyricism of his conversation. And, curiously, having found himself with a Master, he did not believe he had to do what, with that method, is imposed on those initiated: to close his eyes and let them fall. Do you know the game of the cadaver? But they want closed hands and a rigid body: and then, perhaps, \emph{it is foolishness}. ``Parietes non faciunt Christianum'' walls do not make a Christian said Victorinus as a pagan. Then he realized his error and became a Christian.

Another point: I don't like the term ``will''. If by ``will'' one means that which it is meant in the West, let us no longer speak of it: the Orientals do not know it. Will is for Westerners a hard weapon, a stumbling block, a ``construction'': and if one does not mean that, then it is useless to call it ``will'': useless and dangerous, because who wills, in reality, does not will anything: willing is a establishing ties (an example close at hand is in the races of the North of today: the will-obsessed). As to the will-consciousness, it is no longer ``will''.

``Resolution'' is the true term that expresses the super-rational and ascetic reality.

Years ago, down in Africa on an unbearable, infernal, hot, afternoon. I had fallen, naked and immobile, onto the bed. All of a sudden, in the hurricane blast of the sirocco, I felt a dull dark visceral rhythm in the sun. I jumped up from the bed, put on some clothes, and was on the terrace: dazzling sun, sea, clouds of sand, the wind burned the eyes; three black men in front of me, their eyes rolling, the slow powerful rhythm, the bowels twisted inside out, the evasion of the base and the sea, oh! The dark gloomy disheveled sea.

Again, one evening. Jasmine. An Arab courtyard, Arabs lying down, reptilian eyes and a dancing girl, a belly dance and a dance song. The rhythm dug the abyss and the voice filled it up; the rhythm cut through the darkness and the voice called the abyss. And the eyes, oh! The eyes of that woman, the smile of the Virgin, the cosmic dance, the creation-game, the game-rhythm, salient in the voice, the voice bending itself on the rhythm like an abyss on fire.

It is dangerous to approach the \textbf{true} Sufis, those, empty or full, \emph{who have a hidden origin}. The activity will come, if something can still be saved, simply, like passing the self-resolution of chaos in a very clear sky. Like resolving the ties of the seasons. But that will entail anything but certain unclean pseudo-actions of ``willing'', with poor human and also ``super-human'' finalism. The structure asks for, requires absolute disinterest, \emph{the absence of that which makes that which changes, change}.

Atheism consists in \emph{not having a human god} who is concerned with the colds and ailments, a pharmacist god.

To surpass (there are the Chinese), not to bring closer: to act such that heaven overlooks the solid speck: high, very high. Alone then one will see what true immensity is. With the little god close at hand, i.e., close to the soul, because the soul no longer goes far from the hand. Not to open oneself to everyone and so to give a single thing to everyone, but to put oneself very high, in the absolute place, on the lines of a peak like those secure panoramic seats where one glimpses the end of the world, that stays in order to come, of this world and of the other.

Otherwise, let it be said \emph{to men}: ``It is necessary to adore God face down on the ground.''


\flrightit{Posted on 2013-05-12 by Cologero}

\begin{center}* * *\end{center}

\begin{footnotesize}
\begin{sffamily}
\texttt{Logres on 2013-05-13 at 09:59 said:}

``He plays with ghosts: the Lemurian age: individual, humanity, society, family, science. That does not exist. That is outside Tradition: Tradition does not know ghosts: that is why it is not only about showing and suggesting a unique higher type, but of establishing the conditions for which a tradition, which includes all the variously ordered activities but along a single axis, can exist. To hurl thrills of absolute liberation to men who do not even know how basic human existence should be established in the field of action, but always related to a center that is not its essence, is dangerous and absurd. Either life is a rite or it is nothing: either everything reacquires a symbolic character, under the general type of the offering, or else nothing remains.''

He sees straight through ``life'' today. Was it on this site, or somewhere else, that I read that modern individualism will one day be regarded as a species of black magic? Also useful, the endorsement of a species of ``atheism'' \& the definition of ``will'' as resolution. Is it fair to say that the line between Tradition and anti-Tradition is the line between Being and Non-Being? He seems to be saying that taking ``Illusion'' (which I have until understood as ``non-being'') in a certain way can be done along an axis of Being as ``intimation''. On a mere common-sense or basic and vulgar level, could one rephrase this by saying that one's direction is more important than one's current state, unless one is ``awake''?

\hfill

\texttt{Saydan\_Shah on 2015-12-07 at 15:23 said:}

Introduction to Magic


This is the book authored/edited by Evola? And any idea of where I can get more of Guido de Georgio?

\end{sffamily}
\end{footnotesize}

